\documentclass[11pt,a4paper]{article}
\usepackage[utf8]{inputenc}
\usepackage[margin=1in]{geometry}
\usepackage{amsmath,amssymb}
\usepackage{booktabs}
\usepackage{tabularx}
\usepackage{cite}
\usepackage{hyperref}
\usepackage{microtype}

\title{\textbf{Adapting Language Models Under Severe GPU Memory Constraints:\\Architectural Tradeoffs, Quantization Synergy, and Experimental Design}}
\author{\textbf{Research Note}}
\date{\today}

\begin{document}
\sloppy

\maketitle

\begin{abstract}
Adapting large language models under strict hardware limits requires balancing model parameters, optimizer states, gradient graphs, and activations. This research note evaluates six parameter-efficient methods: LoRA, DoRA, AdaLoRA, QLoRA, LoftQ, and QA-LoRA. We separate measured results from design choices, analyze memory trade-offs across scales, provide a structured comparison table, specify a controlled single-GPU experimental protocol, and examine critical methodological boundaries, including an unverified laboratory claim.
\end{abstract}

\section{Architectural Foundations: Decomposition and Dynamic Allocation}

Full fine-tuning of large models is computationally prohibitive under tight memory constraints, as optimizers like Adam require 16 bytes per parameter for gradient and momentum tracking.

Low-Rank Adaptation (LoRA) mitigates this overhead by freezing pre-trained weights $W_0$ and parameterizing updates as $\Delta W = \frac{\alpha}{r} B A$, where $B \in \mathbb{R}^{d \times r}$ and $A \in \mathbb{R}^{r \times k}$ with $r \ll \min(d, k)$ \cite{hu2021lora}. LoRA reduces GPU memory consumption during training by freezing the pre-trained weights and optimizing injected low-rank matrices without storing optimizer states or gradients for the frozen parameters \cite{hu2021lora}. Compared to full fine-tuning of GPT-3 175B, LoRA reduces the number of trainable parameters by 10,000 times \cite{hu2021lora}. On GPT-3 175B, LoRA reduces training VRAM consumption from 1.2TB to 350GB when adapting query and value projection matrices with rank r=4 \cite{hu2021lora}.

While LoRA couples directional and magnitude updates within its low-rank product, Weight-Decomposed Low-Rank Adaptation (DoRA) analyzes their distinct learning dynamics \cite{liu2024dora}. DoRA decomposes pre-trained weights into magnitude and directional components and applies LoRA to update the directional component \cite{liu2024dora}. DoRA enhances both the learning capacity and training stability of LoRA while avoiding additional inference overhead \cite{liu2024dora}. Backpropagation through column-normalized directions introduces intermediate activations. The proposed modification to DoRA leads to a training memory reduction of approximately 24.4\% in fine-tuning LLaMA \cite{liu2024dora}.

Rather than enforcing a static rank across all layers, Adaptive Budget Allocation (AdaLoRA) dynamically shifts parameter capacity toward modules that most influence task loss \cite{zhang2023adalora}. AdaLoRA parameterizes incremental updates in the form of singular value decomposition and prunes redundant singular values based on importance scoring \cite{zhang2023adalora}. To avoid expensive numerical SVD during training, AdaLoRA parameterizes incremental updates to mimic singular value decomposition and adds an orthogonality penalty to regularize singular vectors \cite{zhang2023adalora}. Singular values are pruned along a cubic schedule, concentrating capacity in critical layers.

\section{Quantization Frontiers: Precision, Initialization, and Serving}

Although adapters eliminate optimizer overhead, storing 16-bit base weights remains memory-intensive ($\sim$130\,GB for 65B models).

QLoRA breaks this barrier by compressing the static backbone into 4-bit representation \cite{dettmers2023qlora}. QLoRA reduces memory usage to enable fine-tuning a 65B parameter model on a single 48GB GPU by combining 4-bit NormalFloat quantization, Double Quantization, and Paged Optimizers \cite{dettmers2023qlora}. QLoRA has one low-precision storage data type in 4-bit and one computation data type in 16-bit BrainFloat \cite{dettmers2023qlora}. Weights are dequantized on the fly during computation, while gradients accumulate for 16-bit adapters. Double Quantization compresses the first-stage quantization constants to save an average of approximately 0.37 bits per parameter \cite{dettmers2023qlora}. Paged Optimizers leverage unified memory to mitigate gradient checkpointing memory spikes.

Standard QLoRA quantizes base weights independently while initializing adapters near zero, inducing severe discrepancy in low-bit regimes. LoftQ mitigates this issue by treating low-rank adaptation and quantization as a joint objective \cite{li2023loftq}. LoftQ alternatively applies quantization and low-rank approximation to approximate pre-trained weights, providing an initialization that alleviates quantization discrepancy in LoRA fine-tuning \cite{li2023loftq}. LoftQ initializes the network by jointly optimizing the initial values of the quantized backbone and low-rank adapters \cite{li2023loftq}. Alternating minimization avoids representation collapse. LoftQ achieves 88.0\% accuracy on the MNLI-m dataset using 5 alternating steps \cite{li2023loftq}.

In QLoRA and LoftQ, 16-bit adapters cannot merge directly into 4-bit weights without dequantizing or suffering post-training quantization loss. Quantization-Aware Low-Rank Adaptation (QA-LoRA) resolves this tension by aligning degrees of freedom between quantization and adaptation \cite{xu2023qa}. QA-LoRA uses group-wise operators to balance the degrees of freedom of quantization and adaptation, enabling low-rank adapters to be integrated into quantized weights \cite{xu2023qa}. QA-LoRA enables the quantized base model and auxiliary low-rank weights to be integrated into an integer-quantized model without loss of accuracy \cite{xu2023qa}. This design eliminates runtime dequantization overhead during edge deployment.

\section{Practical Tradeoffs and Comparative Analysis}

Training VRAM decomposes into:
\begin{equation}
M_{\text{total}} = M_{\text{weights}} + M_{\text{optimizer}} + M_{\text{gradients}} + M_{\text{activations}} + M_{\text{workspace}}.
\end{equation}
LoRA eliminates base optimizer states, leaving $M_{\text{weights}}$ as the fixed bound. QLoRA and LoftQ compress $M_{\text{weights}}$ fourfold, making activation memory $M_{\text{activations}}$ (scaling with batch size $B$ and sequence length $L$) the primary bottleneck. While gradient checkpointing curbs activation growth, it incurs $\approx$30\% recomputation overhead. Moreover, on-the-fly dequantization in QLoRA lowers token throughput compared to 16-bit LoRA when memory headroom suffices.

Table~\ref{tab:comparison} explicitly separates architectural design choices from empirically reported measurements, contextualizing each outcome within its specific experimental setting.

\begin{table*}[ht]
\centering
\scriptsize
\caption{Systematic comparison: Architectural characteristics versus measured empirical literature results.}
\label{tab:comparison}
\begin{tabularx}{\textwidth}{lp{2.4cm}p{2.4cm}p{2.9cm}X}
\toprule
& \multicolumn{2}{c}{\textbf{Architectural Design Choices}} & \multicolumn{2}{c}{\textbf{Measured Empirical Results}} \\
\cmidrule(lr){2-3} \cmidrule(lr){4-5}
\textbf{Method} & \textbf{Precision \& Merge} & \textbf{Adapter Formulation} & \textbf{Reported Setting} & \textbf{Measured Literature Outcome} \\
\midrule
LoRA \cite{hu2021lora} & 16-bit base and adapter; lossless FP16 merge & Static rank $\Delta W = \frac{\alpha}{r}BA$ across linear projections & GPT-3 175B, Adam optimizer, $r=4$ (Q, V adapted) & VRAM reduced from 1.2TB to 350GB; 10,000$\times$ fewer trainable params. \\
DoRA \cite{liu2024dora} & 16-bit base and adapter; lossless FP16 merge & Directional LoRA update decomposed from magnitude $m$ & LLaMA-7B commonsense reasoning ($B=16$, accum=4) & Detached norm yields 24.4\% VRAM reduction (37.3\,GB $\to$ 28.2\,GB). \\
AdaLoRA \cite{zhang2023adalora} & 16-bit base and adapter; lossless FP16 merge & Dynamic SVD updates $\Delta W = P\Lambda Q$ with pruning & DeBERTaV3-base on GLUE and natural language generation & Matches full fine-tuning with substantially reduced parameter budget. \\
QLoRA \cite{dettmers2023qlora} & 4-bit NF4 base, 16-bit adapter; dequant. to merge & Static rank LoRA with Double Quantization \& Paged Opt. & LLaMA-65B fine-tuning on single 48GB GPU & Replicates 16-bit full fine-tuning performance; saves $\approx$0.37 bits/param. \\
LoftQ \cite{li2023loftq} & 2-bit or 4-bit base, 16-bit adapter; dequant. to merge & Alternating Frobenius optimization for joint initialization & LLaMA-2 on WikiText-2 and GSM8K; DeBERTa on MNLI & Achieves 88.0\% on MNLI-m with 5 steps; recovers accuracy in $<3$-bit regime. \\
QA-LoRA \cite{xu2023qa} & INT4 base, INT4/FP16 adapter; lossless INT4 merge & Group-wise operators balancing quantization/adaptation degrees & LLaMA-7B/13B and LLaMA-2 on Alpaca \& MMLU benchmarks & Lossless merge into INT4 without PTQ degradation; matches 16-bit LoRA accuracy. \\
\bottomrule
\end{tabularx}
\end{table*}

\section{Proposed Single-GPU Empirical Protocol}

To evaluate these methods without conflating activations with adapter overhead, we specify a controlled benchmark on a 24\,GB GPU (NVIDIA RTX 4090 or A10G) adapting LLaMA-3-8B on Alpaca (52k samples). The primary baseline is 16-bit LoRA ($r=16, \alpha=32$ across all linear projections), evaluated alongside DoRA (detached norm), QLoRA (NF4 base), LoftQ (4-bit initialized), QA-LoRA (INT4 group-wise), and an out-of-memory full fine-tuning reference.

\textbf{Controls:} Random seeds are fixed to 42 across Python, NumPy, and PyTorch CUDA (\texttt{torch.backends.cudnn.deterministic = True}). We fix micro-batch size to $B = 2$, gradient accumulation to 8 (effective batch 16), gradient checkpointing enabled, AdamW ($lr = 2 \times 10^{-4}$), and contexts $L \in \{1024, 2048\}$.

\textbf{VRAM Profiling:} Memory components are isolated via staged PyTorch tracking: (1) static weights $M_{\text{weights}}$ via \texttt{memory\_allocated()} post-loading; (2) static optimizer buffers $M_{\text{opt}}$ after AdamW initialization, giving $M_{\text{static}} = M_{\text{weights}} + M_{\text{opt}}$; (3) dynamic activation memory $M_{\text{act}} = M_{\text{peak\_fwd}} - M_{\text{static}}$ via \texttt{reset\_peak\_memory\_stats()} before the forward pass. The scaling delta $\Delta M_{\text{act}} / \Delta L$ extracts the marginal activation cost per token. Throughput (tokens/s) and MMLU (5-shot) and ARC-Challenge accuracy complete the protocol.

\section{Limitations, Unverified Claims, and Literature Scope}

Rigorous empirical analysis requires acknowledging unverified claims and methodological limits. A pilot method developed in our laboratory reportedly halves training memory compared with LoRA. We designate this statement as strictly unverified: no empirical measurements, hardware details, batching schedules, or activation checkpointing specifications were supplied. Consequently, this claim must not be cited as an established result or adopted as a baseline for experimental planning.

Furthermore, reported empirical gains do not generalize unconditionally. Memory reductions measured during 175B model training with Adam cannot be transferred directly to consumer 7B or 8B model fine-tuning with activation caching under different batch configurations without empirical verification. On smaller models, base weights represent a smaller fraction of peak memory, dampening overall VRAM percentage savings when activation tensors dominate. Additionally, QLoRA's throughput penalty during training dequantization and QA-LoRA's reliance on custom quantized kernels represent practical trade-offs rarely captured by parameter counts alone.

Regarding source accessibility, all six primary papers were registered and verified. We note that automated provider metadata returned a title mismatch warning between OpenAlex and Crossref for the LoRA arXiv entry (arXiv:2106.09685), though full-text verification confirmed identity with the original Low-Rank Adaptation paper.

\bibliographystyle{plain}
\bibliography{refs/references}

\end{document}
